\problemstart{AP-01-011}{NumPy to PyTorch: Shared View Versus Owned Copy}{Difficulty 3/5}{Chapter 1 \textbullet\ pytorch}{Construct two value-equal CPU tensors with deliberately different storage ownership and prove it by mutation.}

        \subsection{Mathematical statement}


For a NumPy array $A$, construct tensors $T_s,T_o$ satisfying
\[
\begin{gathered}
\operatorname{value}(T_s)=\operatorname{value}(T_o)=A,\\
\operatorname{storage}(T_s)\sim\operatorname{storage}(A),\qquad
\operatorname{storage}(T_o)\not\sim\operatorname{storage}(A).
\end{gathered}
\]
The relation $\sim$ is observed through bidirectional mutation, not inferred
from equal values or printed representations.


        \subsection{Code problem}


Using official PyTorch conversion APIs, return a shared CPU tensor and an owned
CPU tensor.  Preserve shape and supported NumPy dtype.  Reject non-writeable or
non-C-contiguous inputs so the educational alias contract remains unambiguous.
Do not move either tensor to a GPU in this problem.


        \begin{contractbox}
        \begin{Code}
        def numpy_to_torch_pair(array: np.ndarray) -> tuple[torch.Tensor, torch.Tensor]:
        \end{Code}
        \end{contractbox}

        \subsection{Theory--engineering gap inventory}

        \begin{gapbox}
        \textbf{Explicit gap.} Equality in a vector space does not encode storage aliasing, device placement, writeability, or which runtime owns the lifetime.

        \textbf{Implicit gap exposed by the result.} The four-way mutation probe distinguishes \code{torch.from\_numpy} from copy construction even though dtype, shape, and initial values match.
        \end{gapbox}

        \subsection{Acceptance evidence}

        \begin{evidencebox}
        \begin{itemize}
          \item NumPy mutation reaches only the shared tensor; shared-tensor mutation reaches NumPy.
  \item Owned-tensor mutation reaches neither NumPy nor the shared tensor.
  \item Both tensors remain on CPU and preserve float32 dtype and shape.
        \end{itemize}
        \end{evidencebox}

        \subsection{Constraints}

        \begin{itemize}
          \item Use \code{torch.from\_numpy} (or documented equivalent) for sharing and an explicit copy API for ownership.
  \item No DLPack, GPU transfer, autograd, or model code.
  \item Install the pinned PyTorch CPU environment to run this problem.
        \end{itemize}

        \begin{sourcebox}
        This is an atomic adaptation, not copied solution code. Nearest primary sources:
        \begin{itemize}
          \item \href{https://docs.pytorch.org/tutorials/beginner/basics/tensorqs_tutorial.html}{PyTorch: Tensors}
  \item \href{https://docs.pytorch.org/docs/stable/generated/torch.from_numpy.html}{PyTorch: torch.from\_numpy}
  \item \href{https://numpy.org/doc/stable/user/basics.copies.html}{NumPy: Copies and views}
        \end{itemize}
        \end{sourcebox}
